\chapter{First- and Second-Order Linear Dynamical Systems}
\label{ch:linsys}
Before any specific circuit, this chapter reviews the linear-systems machinery
that the rest of the book uses. Everything here is generic---it is the same
mathematics whether the system is an RC filter, a mechanical suspension, or a
chemical reactor. Part~III then shows that each amateur-radio circuit is one of
these systems wearing electrical clothing. The payoff of doing it generically
first is that the circuits become instances of a pattern you already know, rather
than a pile of separate formulas.

\section{Four Equivalent Views of a Linear System}
A linear time-invariant system can be written four equivalent ways, and moving
between them is the central skill of this book:
\begin{enumerate}
\item a \textbf{differential equation} relating output \(y(t)\) to input \(u(t)\);
\item a \textbf{state-space model} \(\dot{\state}=A\state+Bu,\; y=C\state+Du\);
\item a \textbf{transfer function} \(G(s)=Y(s)/U(s)\) in the Laplace variable
\(s\);
\item a \textbf{frequency response} \(G(\jj\omega)\), the transfer function on the
imaginary axis.
\end{enumerate}

\begin{controlsbox}
\textbf{Refresher: how the four views connect.} Taking the Laplace transform of
the state-space model (zero initial conditions) and eliminating \(\state\) gives
\[
G(s)=C\,(sI-A)^{-1}B+D .
\]
The \emph{poles} of \(G\) are the eigenvalues of \(A\); they are the system's
natural modes. Setting \(s=\jj\omega\) samples the response to a steady sinusoid
of frequency \(\omega\). So the eigenvalues you compute for stability, the poles
you read on an s-plane, and the peaks you see on a Bode plot are all the same
numbers viewed differently.
\end{controlsbox}

\section{First-Order Systems}
The canonical first-order system has one state and one time constant \(\tau\):
\[
\tau\,\dot{y}+y=K\,u
\qquad\Longleftrightarrow\qquad
\dot{x}=-\frac{1}{\tau}x+\frac{K}{\tau}u,\quad y=x .
\]
Its single pole is real and negative,
\[
G(s)=\frac{K}{\tau s+1},
\qquad
s_{\text{pole}}=-\frac{1}{\tau}.
\]

\subsection{Time-Domain View: the Step Response}
For a unit step input the output rises exponentially to its final value,
\[
y(t)=K\bigl(1-\ee^{-t/\tau}\bigr).
\]
The time constant \(\tau\) is the natural clock of the system:
after one \(\tau\) the response has covered \(1-\ee^{-1}\approx 63.2\%\) of the
way, and after about \(5\tau\) it is settled for practical purposes
(\cref{fig:step1}).

\begin{figure}[htbp]
\centering
\includegraphics[width=0.82\textwidth]{step_first_order}
\caption{First-order step response. One time constant reaches \(63.2\%\); five
time constants are effectively settled.}
\label{fig:step1}
\end{figure}

\subsection{Frequency-Domain View}
On the imaginary axis,
\[
G(\jj\omega)=\frac{K}{1+\jj\omega\tau},
\qquad
|G|=\frac{K}{\sqrt{1+(\omega\tau)^2}},
\qquad
\angle G=-\tan^{-1}(\omega\tau).
\]
The magnitude is flat until the \emph{corner frequency}
\(\omega_c=1/\tau\), then rolls off at \(-\SI{20}{\decibel}\) per decade; the
phase slides from \(0^\circ\) to \(-90^\circ\), passing \(-45^\circ\) at the
corner (\cref{fig:bode1}). The corner frequency is the reciprocal of the time
constant---the same number that governs the transient.

\begin{figure}[htbp]
\centering
\includegraphics[width=0.82\textwidth]{bode_first_order_lowpass}
\caption{First-order low-pass Bode plot. The corner is at \(\omega_c=1/\tau\),
where the gain is \(-\SI{3}{\decibel}\) and the phase is \(-45^\circ\).}
\label{fig:bode1}
\end{figure}

\begin{physicalbox}
One pole tells the whole first-order story three ways: its distance from the
origin \(1/\tau\) is the decay rate in the time domain, the corner frequency in
the frequency domain, and the pole location on the negative real axis in the
s-plane. A first-order system has no imaginary part, so it can never overshoot or
ring.
\end{physicalbox}

\subsection{Worked Example}
A first-order system has \(K=5\) and \(\tau=\SI{2}{\milli\second}\). Its pole is
at \(-1/\tau=-\SI{500}{\per\second}\); the step response reaches \(63.2\%\) of
\(5\) (that is, \(3.16\)) after \SI{2}{\milli\second} and is settled after about
\SI{10}{\milli\second}. The corner frequency is
\(f_c=1/(2\pi\tau)\approx\SI{79.6}{\hertz}\).

\section{Second-Order Systems}
Two independent energy stores give a second-order system, written in
\emph{standard form} as
\[
\ddot{y}+2\zeta\omega_0\,\dot{y}+\omega_0^2\,y=\omega_0^2\,u .
\]
The two parameters are the \textbf{undamped natural frequency} \(\omega_0\) and
the dimensionless \textbf{damping ratio} \(\zeta\). A companion-form state-space
model uses \(x_1=y\), \(x_2=\dot{y}\):
\[
\dot{\state}=
\mat{0 & 1\\ -\omega_0^2 & -2\zeta\omega_0}\state
+\mat{0\\ \omega_0^2}u,
\qquad
y=\mat{1 & 0}\state .
\]
The characteristic polynomial and poles are
\[
s^2+2\zeta\omega_0 s+\omega_0^2=0,
\qquad
s=-\zeta\omega_0\pm\jj\,\omega_0\sqrt{1-\zeta^2}
=-\zeta\omega_0\pm\jj\,\omega_d,
\]
where \(\omega_d=\omega_0\sqrt{1-\zeta^2}\) is the \textbf{damped natural
frequency}.

\begin{controlsbox}
\textbf{Refresher: \(\zeta\), \(Q\), and the damping regimes.} The damping ratio
\(\zeta\) sorts second-order systems into three behaviors:
\begin{description}
\item[\(\zeta>1\) overdamped] two real poles; no oscillation, sluggish approach.
\item[\(\zeta=1\) critically damped] repeated real pole; fastest approach without
overshoot.
\item[\(\zeta<1\) underdamped] complex pole pair; overshoot and ringing.
\end{description}
Radio work usually states the same information as the \emph{quality factor}
\[
Q=\frac{1}{2\zeta},
\]
so a lightly damped, high-\(Q\) resonator and an underdamped control system are
the same object.
\end{controlsbox}

\subsection{Time-Domain View: Damping You Can See}
\Cref{fig:step2} shows step responses across the regimes. A heavily damped
system (\(\zeta\ge1\)) crawls to its final value; a lightly damped one
overshoots and rings for roughly \(Q\) cycles before settling. The peak
overshoot of an underdamped step is
\[
M_p=\exp\!\left(\frac{-\pi\zeta}{\sqrt{1-\zeta^2}}\right).
\]

\begin{figure}[htbp]
\centering
\includegraphics[width=0.82\textwidth]{step_second_order}
\caption{Second-order step responses. Low \(\zeta\) (high \(Q\)) overshoots and
rings; \(\zeta=1\) is the fastest response without overshoot; high \(\zeta\) is
sluggish.}
\label{fig:step2}
\end{figure}

\begin{physicalbox}
\textbf{Mass--spring--damper analogy.} The standard second-order system is
exactly a mass \(m\) on a spring \(k\) with a damper \(c\):
\(m\ddot{x}+c\dot{x}+kx=F\), giving \(\omega_0=\sqrt{k/m}\) and
\(\zeta=c/(2\sqrt{km})\). Underdamped is a plucked guitar string that rings;
overdamped is a screen door with a strong closer that eases shut without
slamming; critically damped is the closer tuned to shut as fast as possible
without bouncing. The RLC circuit of Part~III is this same system with
\(L\leftrightarrow m\), \(1/C\leftrightarrow k\), and \(R\leftrightarrow c\).
\end{physicalbox}

\subsection{The Pole Geometry}
The underdamped poles lie on a circle of radius \(\omega_0\), because
\[
|s|^2=(\zeta\omega_0)^2+\bigl(\omega_0\sqrt{1-\zeta^2}\bigr)^2=\omega_0^2,
\]
and the angle \(\theta\) from the negative real axis satisfies
\(\cos\theta=\zeta\) (\cref{fig:poles2}). The real part \(\zeta\omega_0\) is the
envelope decay rate; the imaginary part \(\omega_d\) is the ringing frequency.
Distance from the origin is \(\omega_0\); angle encodes damping.

\begin{figure}[htbp]
\centering
\includegraphics[width=0.56\textwidth]{poles_second_order}
\caption{Second-order poles for several damping ratios. They ride a circle of
radius \(\omega_0\); the angle from the negative real axis has
\(\cos\theta=\zeta\). \(\zeta=0\) sits on the \(\jj\omega\) axis (undamped).}
\label{fig:poles2}
\end{figure}

\subsection{Frequency-Domain View: the Resonant Peak}
Sampling \(G(\jj\omega)\) near \(\omega_0\) brings the evaluation point close to
the pole, producing a resonant peak whose height grows with \(Q\)
(\cref{fig:bode2}). The flattest response with no peak occurs at
\(Q=1/\sqrt{2}\approx0.707\) (the Butterworth value); above it the magnitude
peaks near \(\omega_0\), and the peak sharpens as the poles approach the
\(\jj\omega\) axis.

\begin{figure}[htbp]
\centering
\includegraphics[width=0.82\textwidth]{bode_second_order}
\caption{Second-order magnitude response. Higher \(Q\) (lower \(\zeta\)) means a
taller, sharper resonant peak, because the poles sit closer to the \(\jj\omega\)
axis.}
\label{fig:bode2}
\end{figure}

\subsection{Worked Example}
A second-order system has \(\omega_0=2\pi(\SI{1}{\kilo\hertz})\) and
\(\zeta=0.1\), so \(Q=1/(2\zeta)=5\). The poles are at
\[
s=-\zeta\omega_0\pm\jj\omega_0\sqrt{1-\zeta^2}
\approx-628\pm\jj\,6250\ \si{\per\second}.
\]
The step overshoot is \(M_p=\exp(-\pi\cdot0.1/\sqrt{1-0.01})\approx0.73\), a
\(73\%\) overshoot, and the resonant peak is about
\(1/(2\zeta\sqrt{1-\zeta^2})\approx5.0\), or \(\SI{14}{\decibel}\)---consistent
with the \(Q=5\) curve in \cref{fig:bode2}.

\section{Stability and a First Look at the Root Locus}
A mode \(\ee^{st}\) decays only if the pole's real part is negative, so a linear
system is stable exactly when \emph{all poles lie in the left half-plane}. Poles
on the \(\jj\omega\) axis give sustained oscillation (the boundary an oscillator
is built to sit on); poles in the right half-plane grow without bound.

When a parameter varies---a resistance, a gain, a feedback fraction---the poles
move, tracing a \emph{root locus}. \Cref{fig:rootlocus2} shows the canonical
second-order locus as the damping increases from zero: the poles start on the
\(\jj\omega\) axis, arc along the circle of radius \(\omega_0\) into the left
half-plane, meet at \(-\omega_0\) at critical damping, then split and travel
apart along the negative real axis. \Cref{ch:splane} develops poles, zeros, and
the root locus in full; the point here is that a single moving pair of poles
captures the entire range from pure oscillation to sluggish overdamping.

\begin{figure}[htbp]
\centering
\includegraphics[width=0.58\textwidth]{rootlocus_second_order}
\caption{Root locus of a second-order system as damping increases. Poles begin
on the \(\jj\omega\) axis (undamped), ride the \(\omega_0\) circle to \(-\omega_0\)
at critical damping, then separate along the real axis (overdamped).}
\label{fig:rootlocus2}
\end{figure}

\begin{exambox}
The exam never asks for \(\zeta\) or a pole plot. It asks about ringing,
bandwidth, selectivity, resonance sharpness, and oscillator stability---every one
of which is a statement about \emph{where the poles are}. Learn the generic
picture here once, and the RC, RL, and RLC circuits of Part~III become three
readings of the same diagram.
\end{exambox}
